%! Sine and Cosine Function Values — Unit 6, Lesson 6.4 — Exit Ticket (Answer Key)
\documentclass[10pt]{article}
\usepackage{precalculus-article}
\usepackage{precalculus-key}   % pulls in -boxes; do NOT also load -boxes
\newcommand{\TallMath}[1]{$\displaystyle #1\rule[-1.4em]{0pt}{3.2em}$}

\begin{document}

\pageheader{Unit 6, Lesson 6.4}{Exit Ticket --- Answer Key}
\namedateperiod

\vspace{0.16in}

\begin{notesbox}{Sine and Cosine Function Values}

\textbf{1.} Fill in the exact values. Show the reference angle and the quadrant you used.

\vspace{0.08in}
\begin{tabularx}{\linewidth}{@{} X X @{}}
$\cos\!\left(\dfrac{2\pi}{3}\right) = $ \ans{$-\dfrac{1}{2}$} &
$\sin\!\left(\dfrac{5\pi}{6}\right) = $ \ans{$\dfrac{1}{2}$} \\[1em]
\end{tabularx}

\textit{Reference angle and quadrant:}

\ans{$2\pi/3$ is in QII; reference angle $= \pi/3$; $\cos(\pi/3) = 1/2$, negate for QII $\Rightarrow -1/2$.\\[3pt]
$5\pi/6$ is in QII; reference angle $= \pi/6$; $\sin(\pi/6) = 1/2$, positive in QII $\Rightarrow 1/2$.}

\vspace{0.16in}

\textbf{2.} $r = 5$, $\theta = 3\pi/4$.

Reference angle $= \pi/4$ (QII). $\cos(\pi/4) = \sqrt{2}/2$, negate for QII.
$\sin(\pi/4) = \sqrt{2}/2$, positive in QII.
\[
  x = 5\cos\!\left(\tfrac{3\pi}{4}\right) = 5\cdot\left(-\tfrac{\sqrt{2}}{2}\right) = -\tfrac{5\sqrt{2}}{2},
  \qquad
  y = 5\sin\!\left(\tfrac{3\pi}{4}\right) = 5\cdot\tfrac{\sqrt{2}}{2} = \tfrac{5\sqrt{2}}{2}
\]

$P = \Bigl($ \ans{$-\dfrac{5\sqrt{2}}{2}$} $,$ \ans{$\dfrac{5\sqrt{2}}{2}$} $\Bigr)$

\end{notesbox}

\end{document}
